\chapter{Deployment Considerations}
\label{ch:deployment}

This chapter collects two deployment-related findings: inference latency
under expected throughput, and an honest comparison against unsupervised
anomaly-detection baselines.

\section{Inference Latency and Throughput}

Inference latency is not the bottleneck for any of the four models. At
batch~=~10{,}000 (a typical NIDS collector burst size), per-flow latency
is 0.11~$\mu$s for logistic regression, 0.80~$\mu$s for XGBoost,
2.80~$\mu$s for LightGBM, and 3.19~$\mu$s for random forest. A
saturated 1~Gbps mirror tap produces approximately 1{,}700 flows per
second --- two orders of magnitude below even the slowest model's
throughput of 313{,}400 flows per second.

\begin{table}[H]
    \centering
    \caption{Per-flow inference latency at batch=10{,}000 on the binary day split.}
    \label{tab:latency}
    \small
    \begin{tabular}{lcc}
        \toprule
        \textbf{Model} & \textbf{$\mu$s / flow (bs=10k)} & \textbf{Flows / second} \\
        \midrule
        LogReg & 0.11 & 9{,}108{,}090 \\
        XGBoost & 0.80 & 1{,}248{,}123 \\
        LightGBM & 2.80 & 356{,}793 \\
        Random Forest & 3.19 & 313{,}400 \\
        \bottomrule
    \end{tabular}
\end{table}

Model selection should therefore be driven by accuracy, robustness, and
operational simplicity, not by latency. Any of the four models can keep
up with a typical enterprise mirror tap on a single CPU core; latency
becomes a factor only at much higher throughputs (saturated 100~Gbps
backbones), which are out of scope here.

\section{Unsupervised Anomaly Detection as an Honest Baseline}

A natural question is whether the labels are doing real work --- could
the same detection performance be achieved without them? To answer, two
unsupervised baselines were trained on Monday-to-Wednesday benign flows
only: Isolation Forest and Local Outlier Factor (LOF). Both flag Friday
anomalies without ever having seen a labelled attack.

Isolation Forest reaches ROC-AUC $= 0.807$ on the binary day task; LOF
reaches $0.849$. Both are substantially worse than the supervised
random forest at ROC-AUC $= 0.940$, which confirms that the attack
labels carry real signal. The gap widens at strict FPR budgets: at 1\%
FPR, Isolation Forest catches only about 1\% of attacks while the
supervised random forest catches 63\%. Unsupervised detection is not
bad in the abstract --- it works without labels, which is valuable ---
but at the FPR budgets where a real SOC operates, supervised wins
decisively. The practical recommendation is to use unsupervised
detection as a novelty layer alongside the supervised detector for
known attack families, rather than as a replacement.
